\section{Reflection Perspective}\label{reflection-perspective}

This section will focus on our reflections about the process, some of
our major issues and what we have learned for improving the process in
the future.

\subsection{Devops-style work}
This project have shown how powerfull it can be to bridge the gap between developer and operations.
By automating processes, such as a full CI/CD pipeline, the distance and complexity of going from a feature branch to development have become small and manageable.
The project created a short feedback loop, which helped opening up for more continous work, both in the world of development, Devops and operations.
To help enable a short feedback loop, we have tried to keep our work visible. 
This was done though the use of conventional commits, a notebook for changes, a project board, and group chat for emergencies.
This also allowed for showing changes and thereby fueling continual learning.
An example of an issue in our project board can be seen here \url{https://github.com/DevBobs-devops/Devbobs_MiniTwit/pull/66}. \\

As the project grew it became harder to diagnose and understand why something failed.
With the monitoring and logging setup this became easier, as everything was collected in one place and allowed for quicker feedback. 
In retrospective our setup could have been setup better, which mostly comes down to setting up the dashboards better.
We will touch more on this in \ref{reflection-logging}.



\subsection{Biggest challenges during
development}\label{biggest-issues-during-development}
This section will touch some of the biggest issues we had with regards to evolution, and refactoring and operation, and maintenance.
\subsubsection{Migrating Database}\label{migrating-database}
The project started by using our old Chirp project with a sqlite database. 
This gave us some difficulties when migrating database, as data in the old database did not fit the new Postgres database. 
We solved this by converting the sqlite dump to csv-files, which could then be converted into postgres dumps in the correct format. 
In reflection using Postgres from the beginning would have been better if we knew the project would grow big, as the Sqlite database would create technical debt we later had to pay for.
This we could not have know at the time we built Chirp, but shows how evolution of a system will require refactoring later if such evolution was not thought of when building the system.
We also updated our old source code to adapt to the new changes when relevant, attempting to remove possible errors caused by our old source code.


\subsubsection{Logging}\label{reflection-logging}

Logging has been a constant issue. Logs would ``dissapear'' after a
while using Grafana log drilldown, only showing logs for one or two of
our containers. We instead made a log dashboard to show the logs, like the minitwit
dashboard, fixing the problem. 

Having more indept logging, would have made it easier to spot errors in
our program. We often had to check the logs on the server, by ssh into
it, instead of using the grafana dashboard. Having more logs show up in
grafana, for example more statistics about queries to the database, could
have given us more information about the more niche parts of our program
and errors.

\subsubsection{Test-server}\label{test-server}

Having a test-server would have been benefitial for this project, as we
currently only have our ``main'' server. This would have allowed us to
further test new features or fixes, before they ended up in production.
This also holds true for the database, as we often had to be very
cautious out of fear that we would overwrite or delete data. We did make
another DigitalOcean project near the end of the project, to test Terraform and docker
swarm, however this never got added to our pipeline.

\subsubsection{Updating test-file}\label{updating-test-file}

Our current test run on an old docker compose file, from before we
introduced terraform and docker swarm. It would be better
practice to use the stack file for tests, rather than the compose
file, for more reflective test results. This highlights the importance of active
updating older services, to remove outdated legacy code. 

\subsubsection{Notebook}\label{notebook}

We implemented a markdown file, notebook.md, to document the progress of 
the project, including any problems we encountered during a task and how we 
solved it. It was used to facilitate "organisational learning" in the group, 
and became the primary mechanism for how group members shared local 
discoveries with the rest. 